% ============================================================
% Section 9: Composite Loss
% ============================================================

\section{Composite Loss}
\label{sec:loss}

\aletheion{} uses a composite loss with 14 components, grouped into
standard cross-entropy and 13 epistemic losses with progressive
annealing.

\subsection{General Formulation}

\begin{equation}
    \calL_{\text{total}} = \lambda_{\text{CE}} \cdot \calL_{\text{CE}} + \alpha(t) \cdot \sum_{i=1}^{13} \lambda_i \cdot \calL_i
    \label{eq:total_loss}
\end{equation}

where $\alpha(t) \in [0, 1]$ is the annealing factor and $\lambda_i$
are per-component weights.

\subsection{Cross-Entropy}

The primary loss is the standard cross-entropy for autoregressive
language modeling:

\begin{equation}
    \calL_{\text{CE}} = -\frac{1}{|\calS|} \sum_{t \in \calS} \log P(x_t \mid x_{<t})
\end{equation}

where $P(x_t \mid x_{<t}) = \text{softmax}(\bm{l}_t / \tau)_{x_t}$ and
$\calS$ excludes padding tokens.

\subsection{12 Epistemic Losses}

\begin{table}[H]
\centering
\caption{Composite loss components with default weights.}
\label{tab:losses}
\begin{tabular}{clcp{6cm}}
\toprule
\textbf{\#} & \textbf{Name} & \textbf{$\lambda$} & \textbf{Objective} \\
\midrule
1 & VARO & 0.10 & Calibrate $q_1 \cdot q_2$ vs accuracy \\
2 & VI & 0.01 & Maintain $\phi \geq \phi_{\text{critical}}$ + severity penalty \\
3 & MAD Calibration & 0.05 & Calibrate confidence vs accuracy \\
4 & Metric & 0.001 & $\kappa(\bG) \to 1$ (condition number) \\
5 & Eidos & 0.005 & Axis balancing \\
6 & Conflict & 0.005 & Minimize $\gamma^2$ \\
7 & Consciousness & 0.003 & Energy $\geq 0.2$ \\
8 & Grounding & 0.005 & Task entropy + calibrate ambiguity \\
9 & Plasticity & 0.002 & Plasticity $\geq 0.3$ \\
10 & Frontier & 0.002 & Maximize exploration \\
11 & MOPsi & 0.003 & $\psi$-confidence alignment \\
12 & Contrastive & 0.003 & Anti-collapse of views \\
\bottomrule
\end{tabular}
\end{table}

\subsection{Loss Details}

\paragraph{VARO (Calibration).}
Calibrates the combined uncertainty signal against actual correctness:

\begin{equation}
    \calL_{\text{VARO}} = \E\big[(q_1 \cdot q_2 - p_{\text{correct}})^2\big]
\end{equation}

where $p_{\text{correct}} = P(x_t = x_t^* \mid x_{<t})$ is the softmax
probability of the correct token. This forces the product $q_1 \cdot q_2$
to estimate the model's actual confidence in each prediction.

\paragraph{VI Regularization.}
Penalizes degraded manifold and excessive severity (\Cref{sec:vi}):

\begin{equation}
    \calL_{\text{VI}} = \E\big[\text{ReLU}(\phi_{\text{critical}} - \phi)^2\big] + 0.3 \cdot \E[s^2], \quad \phi_{\text{critical}} = 0.5
\end{equation}

\paragraph{MAD Calibration.}
BCE between confidence and accuracy (\Cref{eq:mad_loss}).

\paragraph{Metric Regularization.}
\begin{equation}
    \calL_{\text{metric}} = \log(\kappa(\bG) + 1)
\end{equation}

\paragraph{Eidos.}
Penalizes axis imbalance:

\begin{equation}
    \calL_{\text{eidos}} = \E\big[(\text{std}(\text{axis\_balance}) - 0.15)^2\big]
\end{equation}

\paragraph{Conflict.}
Minimizes $\phi$-$\psi$ conflict intensity:

\begin{equation}
    \calL_{\text{conflict}} = \E[\gamma^2]
\end{equation}

\paragraph{Consciousness.}
Prevents energy collapse:

\begin{equation}
    \calL_{\text{cons}} = \E\big[\text{ReLU}(0.2 - \text{energy})^2\big]
\end{equation}

\paragraph{Grounding.}
Task probability entropy + ambiguity calibration:

\begin{equation}
    \calL_{\text{ground}} = -\E\Big[\sum_k p_k \log p_k\Big] + \E\big[(\text{ambiguity} - q_1)^2\big]
\end{equation}

\paragraph{Plasticity.}
Prevents total depletion:

\begin{equation}
    \calL_{\text{plast}} = \E\big[\text{ReLU}(0.3 - \text{plast})^2\big]
\end{equation}

\paragraph{Frontier.}
Maximizes exploration of novel regions:

\begin{equation}
    \calL_{\text{frontier}} = -\E\big[f_{\text{frontier}} \cdot \log(\text{novelty} + \epsilon)\big]
\end{equation}

\paragraph{MOPsi.}
Aligns $\psi$ with confidence:

\begin{equation}
    \calL_{\text{mopsi}} = \E\big[(\psi_{\text{score}} - \kappa)^2\big]
\end{equation}

\paragraph{Contrastive.}
Anti-collapse: ensures the two views are not identical:

\begin{equation}
    \calL_{\text{contr}} = -\E[\log(\text{div} + \epsilon)] + 0.1 \cdot \E\big[\text{ReLU}(\text{div} - 0.8)^2\big]
\end{equation}

The first term maximizes divergence (anti-collapse), the second prevents
excessive divergence.

\subsection{Annealing Schedule}
\label{sec:annealing}

Epistemic losses are introduced gradually to stabilize training:

\begin{equation}
    \alpha(t) = \begin{cases}
        0 & \text{if } t < 0.1 \cdot T_{\text{total}} \\
        \frac{t - 0.1T}{0.4T} & \text{if } 0.1T \leq t < 0.5T \\
        1 & \text{if } t \geq 0.5T
    \end{cases}
    \label{eq:annealing}
\end{equation}

\begin{itemize}
    \item \textbf{Phase 1} (0--10\%): Only $\calL_{\text{CE}}$. The model
    learns to distribute tokens before any epistemic signal.

    \item \textbf{Phase 2} (10--50\%): $\alpha$ grows linearly from 0 to 1.
    Epistemic losses are smoothly introduced.

    \item \textbf{Phase 3} (50--100\%): $\alpha = 1$. All losses active
    at full weight.
\end{itemize}

\begin{theorem}[Annealing Stability]
The annealing schedule ensures that epistemic gradients do not dominate
CE gradients in early phases:

\begin{equation}
    \frac{\|\nabla_\theta \alpha \calL_{\text{epist}}\|}{\|\nabla_\theta \calL_{\text{CE}}\|} \leq \alpha(t) \cdot \frac{\sum_i \lambda_i \|\nabla_\theta \calL_i\|}{\|\nabla_\theta \calL_{\text{CE}}\|}
\end{equation}

For $\alpha(t) \ll 1$ in the initial phases, the ratio is controlled.
\end{theorem}

\subsection{Component Weights}

The weights $\lambda_i$ were calibrated empirically using the following
criteria:

\begin{enumerate}
    \item $\lambda_{\text{CE}} = 1.0$ (anchor)
    \item Regularization losses ($\leq 0.01$): prevent degeneration
    without dominating
    \item Calibration losses ($\leq 0.005$): fine-tune signals
    \item Exploration losses ($\leq 0.003$): encourage diversity
\end{enumerate}
